\documentclass[11pt]{article}
\usepackage[margin=1.15in]{geometry}
\usepackage{amsmath, amssymb, amsthm}
\usepackage{booktabs}
\usepackage{xcolor}
\definecolor{linknavy}{RGB}{25, 55, 125}
\usepackage[colorlinks=true, linkcolor=linknavy, citecolor=linknavy, urlcolor=linknavy]{hyperref}
\usepackage{orcidlink}
\usepackage{fancyhdr}
\newtheorem{corollary}{Corollary}
\newcommand{\C}{\mathbb{C}}
\newcommand{\Q}{\mathbb{Q}}

% Page-1 header block (conventions/manuscript-header-standard.md). The four document
% macros are set by the publishing tool when the version DOI is reserved. At this text
% width the two DOIs do not fit on one line of the box, so each has its own line.
\newcommand{\doctype}{Preprint}
\newcommand{\docversion}{v0.06}
\newcommand{\docdate}{18 September 2026}
\newcommand{\versiondoi}{10.5281/zenodo.22834179}
\newcommand{\conceptdoi}{10.5281/zenodo.21503371}
\newcommand{\docrepo}{github.com/fsantibanezleal/CAOS\_RESEARCH}

\newcommand{\headerblock}{%
\begin{center}
\fbox{\parbox{0.94\textwidth}{\small\raggedright
\textbf{\doctype} \quad$\cdot$\quad \docversion, \docdate \quad$\cdot$\quad Licence CC BY 4.0\\[2pt]
CAOS Research open-problems programme \quad$\cdot$\quad problem \texttt{jacobian-conjecture}\\[2pt]
Version DOI \href{https://doi.org/\versiondoi}{\texttt{\versiondoi}}\\[2pt]
Concept DOI (always latest) \href{https://doi.org/\conceptdoi}{\texttt{\conceptdoi}}\\[2pt]
Not peer reviewed. Source code, data and computational records: \href{https://\docrepo}{\texttt{\docrepo}}
}}
\end{center}
\vspace{4pt}}

\title{The Collateral Impact of the\\ Three-Dimensional Jacobian Counterexample:\\ a Source-Verified Cascade}
\author{Felipe Santib\'a\~nez-Leal\,\orcidlink{0000-0002-0150-3246}\\
\small CAOS Research program, \texttt{github.com/fsantibanezleal/CAOS\_RESEARCH}}
\date{\docdate}

\pagestyle{fancy}
\fancyhf{}
\fancyhead[L]{\small CAOS Research \doctype}
\fancyhead[R]{\small The Consequence Cascade \textperiodcentered\ \docversion}
\fancyfoot[C]{\small \thepage}
\renewcommand{\headrulewidth}{0.4pt}
\fancypagestyle{plain}{\fancyhf{}\fancyfoot[C]{\small \thepage}\renewcommand{\headrulewidth}{0pt}}

\begin{document}
\maketitle
\headerblock

\begin{abstract}
On 2026-07-20, Levent Alp\"oge announced an explicit counterexample to the Jacobian conjecture
in dimension 3, found working with the language model Claude Fable 5 (Anthropic). The credit for % voice-ok: external attribution, as the announcement states it
every consequence documented here belongs to that single map. This paper documents what else
that map decides: through implication chains published over five decades, each re-verified
against its primary source (EXP-001, EXP-016, EXP-018, \mbox{EXP-136}), the counterexample resolves or
partially resolves the Mathieu conjecture, the Dixmier conjecture, the Poisson conjecture, the
Gaussian moments conjecture, Zhao's vanishing conjecture, the Image conjecture, and the
symmetric (gradient) form of the Jacobian conjecture. We state each corollary with its chain
and its reference, list precisely what remains open, and describe the verification method,
which turns a one-line certificate into a status map of a whole neighborhood of problems. The
symmetric/gradient and vanishing corollaries are made explicit: a non-invertible gradient
Keller map $\mathrm{id} + \nabla f$ on $\C^{48}$, with $f$ a homogeneous Hessian-nilpotent
quartic (382 monomials over $\Q(i)$) and an exact two-point collision, built from Thompson's
dimension-24 cubic-homogeneous form by the de Bondt--van den Essen construction and verified
end to end in exact arithmetic.
\end{abstract}

\section{The counterexample and its verification}

The counterexample is the polynomial map $F : \C^3 \to \C^3$, $u = 1 + xy$,
\[
F = \bigl(u^3 z + y^2 u (4 + 3xy),\;\; y + 3x u^2 z + 3x y^2 (4 + 3xy),\;\; 2x - 3x^2 y - x^3 z\bigr),
\]
with $\det JF = -2$ everywhere and
$F(0, 0, -\tfrac14) = F(1, -\tfrac32, \tfrac{13}{2}) = F(-1, \tfrac32, \tfrac{13}{2})$
\cite{Alpoge2026}. We validated it independently in exact rational arithmetic and computed the
full fiber over the repeated value: it consists of exactly the three announced points
(EXP-001). Every consequence below is a corollary of this validation plus one or more published
implications; for each implication the primary source was located and its direction and
dimension index confirmed (EXP-016, EXP-018 and EXP-136). The corollaries of
Section~\ref{sec:cascade} add no mathematics beyond this verification and the composition of the
published implications, except the statement that $F$ is not affinely of homogeneous type,
checked by a Gr\"obner basis computation (EXP-136); the explicit witness of
Section~\ref{sec:witness} is a construction from published ingredients (Thompson's
cubic-homogeneous form and the de Bondt--van den Essen reduction).

\section{The cascade}\label{sec:cascade}

\begin{corollary}[Smale's 16th problem]
The Jacobian conjecture is false for every $N \ge 3$; the two-variable case remains open.
\end{corollary}

\begin{corollary}[Mathieu conjecture, 1997]
Mathieu proved that his conjecture, taken for all compact connected Lie groups, implies the
Jacobian conjecture \cite{Mathieu1997}; Duistermaat and van der Kallen proved the abelian case
\cite{DvdK1998}. At a fixed $N$ the argument shows that the conjecture for $SU(N)$ makes every
Keller map $x - h$ of $\C^N$ with $h$ homogeneous invertible
\cite[Theorems~4.16 and~4.23]{Zwart2025}. The cubic-homogeneous form of the counterexample on
$\C^{24}$ (Section~\ref{sec:witness}) is such a map and is not injective, so the Mathieu
conjecture is \textbf{false for $SU(N)$, $N \ge 24$}. For $SU(2)$ an explicit counterexample is
given in \cite{LongSU2}. For $3 \le N \le 23$ the argument does not reach $F$: its components
have total degrees $7, 6, 4$ and its second partial derivatives have no common zero, so $F$ is
not affinely equivalent to a map $x - h$ with $h$ homogeneous (EXP-136), and the reduction of
$F$ to that form adds variables \cite{BCW1982}.
\end{corollary}

\begin{corollary}[Dixmier conjecture, 1968]
$JC(2n)$ implies $\mathrm{Dixmier}(n)$, and $\mathrm{Dixmier}(n)$ implies $JC(n)$
\cite{BKK2007, Tsuchimoto2005}. Since $JC(n)$ is false for every $n \ge 3$, the second
implication gives that $\mathrm{Dixmier}(n)$ is false for every $n \ge 3$: the
\textbf{full Dixmier conjecture is false}. Ranks $1$ (Dixmier's original question) and $2$
remain open; the first implication gives no information on rank $2$, because it runs from
$JC(4)$ to $\mathrm{Dixmier}(2)$. A proof of rank $1$ is claimed in \cite{Zheglov2024}.
\end{corollary}

\begin{corollary}[Poisson conjecture, 2007]
For each $n$, $JC(2n) \Rightarrow \mathrm{Poisson}(n) \Rightarrow \mathrm{Dixmier}(n)
\Rightarrow JC(n)$ in any characteristic, where $\mathrm{Poisson}(n)$ concerns the canonical
Poisson algebra in $2n$ variables \cite[Theorem~7]{AvdE2007}. Hence $\mathrm{Poisson}(n)$ is
false for every $n \ge 3$: the \textbf{full Poisson conjecture is false}. Indices $1$ and $2$
remain open, for the same reason as the Dixmier ranks.
\end{corollary}

\begin{corollary}[Gaussian moments conjecture]
GMC implies JC \cite{DEZ2017}; hence \textbf{GMC is false} at some finite dimension.
\end{corollary}
\noindent An explicit counterexample in three variables, with $P$ of five terms and total
degree $4$, is given in \cite{Long2026}; GMC is therefore false in every dimension $n \ge 3$.
GMC holds in dimension $1$ \cite[Proposition~4.2]{DEZ2017}, and a proof of the two-variable case
is claimed in \cite{Wilson2026}.

\begin{corollary}[Zhao's vanishing conjecture]
The vanishing conjecture for Hessian-nilpotent polynomials is equivalent to JC
\cite{Zhao2004, Zhao2007}; hence it is \textbf{false}: there exists a Hessian-nilpotent
quartic $P$ with $\Delta^m P^m = 0$ for all $m$ but $\Delta^m P^{m+1} \ne 0$ for infinitely
many $m$. \emph{An explicit witness is given in Section~\ref{sec:witness}.}
\end{corollary}

\begin{corollary}[Image conjecture]
The Image conjecture implies the vanishing conjecture \cite{vdE2010}; hence it is
\textbf{false in some dimension}.
\end{corollary}

\begin{corollary}[Symmetric / gradient Jacobian conjecture]
JC reduces to cubic homogeneous maps with symmetric Jacobian \cite{dBvdE2005}; with the Zhao
equivalence this yields a dimension in which a \textbf{gradient Keller map is
non-invertible}. \emph{An explicit such map, in dimension 48, is given in
Section~\ref{sec:witness}.}
\end{corollary}

\section{The explicit witness in dimension 48 (EXP-041)}\label{sec:witness}

As of 2026-07-22 no explicit symmetric, gradient, or Hessian-nilpotent witness had been
published (on a search of the literature; the July 2026 consequence notes are non-constructive
in dimension). We construct one in three steps, each verified in exact arithmetic.

\emph{Input.} W. Thompson's explicit Bass--Connell--Wright cubic-homogeneous form of the
counterexample \cite{Thompson2026}: $G = U + H$ on $\Q^{24}$, $H$ cubic homogeneous with 54
monomials, with a published rational collision pair. We verified this form independently
(homogeneity; the exact collision; nilpotency of $JH$ by an exact sparse chain). The
nilpotency index is $18$, not $17$ as stated there: $(JH)^{18} = 0$ while $(JH)^{17}$ retains
six terms.

\emph{Construction} (de Bondt--van den Essen \cite{dBvdE2005}). The potential
$f = -i \sum_{j=1}^{24} H_j(x + iy)\, y_j$ is a homogeneous quartic in 48 variables (382
monomials over $\Q(i)$) with nilpotent Hessian (their characteristic-polynomial identity
reduces this to the nilpotency of $JH$). The triangular conjugation
$S^{-1} \circ (\mathrm{id} + \nabla f) \circ S = (x + H(x),\; y + JH(x)^t y - iH(x))$ was
verified identically in all 48 components, so
$\widetilde F = \mathrm{id} + \nabla f$ is a symmetric Keller map ($\det J\widetilde F = 1$).

\emph{The collision.} With $(u, v)$ Thompson's collision pair and
$w = i\,(I + JH(v)^t)^{-1}(u - v)$, the points $p_1 = (u, 0)$ and $p_2 = (v - iw,\, w)$ are
distinct and satisfy $\widetilde F(p_1) = \widetilde F(p_2)$ exactly over $\Q(i)$.
\textbf{Hence the Hessian conjecture fails in dimension 48 with an explicit witness}, and
$P_\star = -f$ is an explicit homogeneous Hessian-nilpotent quartic falsifying the vanishing
conjecture (the finite certificate of failure is the collision itself; the consistency check
$\Delta P_\star^2 \ne 0$, forced by Zhao's own Corollary 3.9 for any failing potential,
holds with 9028 terms). The reality obstruction is intrinsic: a nonzero real
Hessian-nilpotent polynomial does not exist, so the $\Q(i)$ coefficients are unavoidable.
Symmetric positive results (dimension $\le 4$ nonhomogeneous, $\le 5$ homogeneous, de
Bondt--van den Essen and van den Essen--Washburn) leave the minimal witness dimension open
between those bounds and 48.

\section{What remains open}

The two-variable Jacobian conjecture; Dixmier and Poisson at ranks 1 and 2 (a proof of
Dixmier's rank 1 is claimed in \cite{Zheglov2024}); Mathieu for $SU(N)$ with $3 \le N \le 23$
and for the other non-abelian compact groups; the moments statement in dimension $2$ (a proof is
claimed in \cite{Wilson2026}); the minimal failing
dimensions of the vanishing and Image statements (for the symmetric/vanishing pair the
witness of Section~\ref{sec:witness} pins the failing dimension at $\le 48$, with the known
positive results leaving the minimum open above $5$); and explicit witnesses for the remaining
non-constructive rows. Not included in this list: the Markus--Yamabe conjecture (refuted
independently in 1997, not a corollary of the counterexample) and statements whose implication
direction does not propagate the falsity of JC.

\section{Method and data availability}

{\raggedright
All computations use exact arithmetic (no floating point) and are archived in the repository
\url{https://github.com/fsantibanezleal/CAOS_RESEARCH}. An identifier EXP-$nnn$ denotes one
archived record in the experiment folders of the problem,
\path{problems/algebraic-geometry/jacobian-conjecture/experiments/}: a computation with its
script and output, or a check of published implications against their primary sources; each
record states its question, method and result. The records used here are EXP-001 (exact
validation, full fiber), EXP-016 and EXP-018 (each implication located in its primary source and
its direction confirmed), EXP-041 (the dimension-48 witness, with its input data, script and
outputs) and EXP-136 (the dimension index of each implication; the homogeneous-type test of $F$;
exact checks of the explicit counterexamples of \cite{Long2026, LongSU2} for the first powers).\par}

\begin{thebibliography}{99}\raggedright
\bibitem{Alpoge2026} L. Alp\"oge, announcement of the counterexample, X (@\_\_alpoge\_\_),
2026-07-20; found working with Claude Fable 5 (Anthropic); question posed by Akhil. % voice-ok: external attribution, as the announcement states it
\bibitem{Mathieu1997} O. Mathieu, Some conjectures about invariant theory and their
applications, in \emph{Alg\`ebre non commutative, groupes quantiques et invariants} (Reims,
1995), S\'emin. Congr. 2, Soc. Math. France, 1997, 263--279.
\bibitem{LongSU2} C. D. Long, Counterexamples to the $xz$-conjecture and the Mathieu conjecture
for $SU(2)$, arXiv:2607.19012, 2026.
\bibitem{BCW1982} H. Bass, E. H. Connell, D. Wright, The Jacobian conjecture: reduction of
degree and formal expansion of the inverse, \emph{Bull. Amer. Math. Soc. (N.S.)} 7 (1982),
287--330.
\bibitem{DvdK1998} J. J. Duistermaat, W. van der Kallen, Constant terms in powers of a Laurent
polynomial, \emph{Indag. Math.} 9 (1998), 221--231.
\bibitem{Zwart2025} K. Zwart, Mathieu's approach to the Jacobian Conjecture (review),
arXiv:2511.16561.
\bibitem{BKK2007} A. Belov-Kanel, M. Kontsevich, The Jacobian conjecture is stably equivalent
to the Dixmier conjecture, arXiv:math/0512171; \emph{Mosc. Math. J.} 7 (2007).
\bibitem{Tsuchimoto2005} Y. Tsuchimoto, Endomorphisms of Weyl algebra and $p$-curvatures,
\emph{Osaka J. Math.} 42 (2005), 435--452.
\bibitem{Zheglov2024} A. Zheglov, The conjecture of Dixmier for the first Weyl algebra is true,
arXiv:2410.06959, 2024 (revised 2026).
\bibitem{AvdE2007} P. K. Adjamagbo, A. van den Essen, On the equivalence of the Jacobian,
Dixmier and Poisson conjectures in any characteristic, arXiv:math/0608009; \emph{Acta Math.
Vietnam.} 32 (2007), 209--218.
\bibitem{DEZ2017} H. Derksen, A. van den Essen, W. Zhao, The Gaussian moments conjecture and
the Jacobian conjecture, \emph{Israel J. Math.} 219 (2017), 917--928; arXiv:1506.05192.
\bibitem{Long2026} C. D. Long, Small counterexamples to the Gaussian moments conjecture,
arXiv:2607.18186, 2026.
\bibitem{Wilson2026} M. Wilson, A face-isolation proof of the two-variable Gaussian moments
conjecture, arXiv:2607.23887, 2026.
\bibitem{Zhao2004} W. Zhao, Hessian nilpotent polynomials and the Jacobian conjecture,
arXiv:math/0409534; \emph{Trans. Amer. Math. Soc.} 359 (2007).
\bibitem{Zhao2007} W. Zhao, A vanishing conjecture on differential operators with constant
coefficients, arXiv:0704.1691.
\bibitem{vdE2010} A. van den Essen, The amazing Image conjecture, arXiv:1006.5801.
\bibitem{dBvdE2005} M. de Bondt, A. van den Essen, A reduction of the Jacobian conjecture to
the symmetric case, \emph{Proc. Amer. Math. Soc.} 133 (2005), 2201--2205.
\bibitem{Thompson2026} W. Thompson, Explicit cubic-homogeneous Jacobian counterexample
(Bass--Connell--Wright form of the Alp\"oge map, dimension 24), 2026-07-21;
\path{github.com/wtho704/explicit-cubic-homogeneous-jacobian-counterexample}; Zenodo
record 21466221.
\end{thebibliography}

\end{document}
